Jõudlustestimise eesmärgiks on saada teada kui kiire ja vastupidav on mingi tarkvaralahendus.
Selle lõputöö käigus viiakse läbi jõudlustestimine loodud näidislahenduse peal. 
Testimise skoopi kuuluvad ka Harmony SMP ning Phoss SMP server.

Jõudlustestimise skoopi ei kuulu WS plugin sõnumite saatmine läbi juurdepääsupunkti. 
Testimise eesmärgiks on teha kindlaks loodud näidislahenduse tulemused võrreldes teiste komponentidega.
Sellest tulenevalt ei ole vajalik testida Harmony juurdepääsupunkti jõudlust.

Testimine jaguneb järgnevateks etappideks:

\begin{enumerate}
  \item Näidislahenduse SMP ja SML testimine
  \item Phoss SMP testimine
  \item Harmony SMP testimine
\end{enumerate}


\section{Metoodika}
Jõudluse testimisel tuleb paika seada millised on testimise sisendid ning
millised on võrreldavad väljundid.

Väljundina võrreldakse töödeldud päringute kogust sekundis ning päringute 99\% latentsust.
Seda tüüpi tulemuste kogumiseks on kõige paremaks tööriistaks Apache Benchmark. Apache Benchmark (ab)
on lihtne, kuid võimas käsureatööriist, mida kasutatakse veebiserverite (mitte ainult Apache) 
koormustestimiseks ja jõudluse mõõtmiseks. See annab kiire ülevaate sellest, kui hästi suudab 
server tulla toime suure hulga päringutega. \cite{apache_ab_tool}

Apache Benchmark tööriista puhul peab määrama mitu päringut tehakse kokku ning kui palju
päringuid tehakse samaaaegselt. Autor tegi jõudlustestimiseks ab argumentide ja tulemuste malli
(tabel~\ref{tab:ab-performance-table-template}). Päringute arv ja samaaegsus on järjest kasvav.
Nii on võimalik võrrelda tarkvara jõudlust erinevatel koormusastmetel.

Apache Benchmark moodi lahendus DNS-i maastikul on \texttt{dnsperf}. \texttt{dnsperf} on tasuta tööriist,
mille on välja töötanud Nominum/Akamai (2006–2018) ja DNS-OARC (alates 2019. aastast)
ning mis teeb latentsuse ja läbilaskevõime meetrikate kogumise lihtsaks \cite{dnsperf_tool}.

\begin{table}[H]
    \caption{ab jõudlustestimise mall}
    \label{tab:ab-performance-table-template}
    \begin{tabular}{|c|c|c|c|}
        \hline
        Päringute arv (-n) & Samaaegsed päringud (-c) & Päringuid sekundis (mediaan) & 99\% latentsus \\
        \hline
        100 & 10 &  &  \\
        \hline
        1000 & 100 &  &  \\
        \hline
        5000 & 250 &  &  \\
        \hline
        10000 & 1000 &  &  \\
        \hline
    \end{tabular}
\end{table}

Kõik testimised toimuvad kasutades Intel(R) Core(TM) Ultra 7 255H (16) @ 5.10 GHz protsessorit.


\section{Näidislahenduse SMP}
Autor käivitas näidislahenduse SMP/SML serveri. 
Apache Benchmark tulemused